\subsection{辐射反作用的应用}

前面建立了辐射反作用的两条主线：在无反作用轨道上进行频谱分析，以及由 Dirac 方程降阶得到的缓慢轨道反作用。本节把两条主线合并起来，应用到两个具体体系：辐射振子，以及有限质量二体库仑运动。后者先沿无反作用轨道求完整频谱，再把频谱积分给出的能量、角动量通量绝热地反馈到轨道常数。

\subsubsection{辐射振子、自然线宽与散射截面}

用辐射振子说明同一个 $\mathsf M(\omega)$ 如何同时决定附加质量、频移、线宽和散射。设恢复力为 $-K\vb x$，在尚未取点极限时，由前面自作用总有效质量的频域运动方程直接得到
\begin{equation*}
	\left[K-\omega^2\mathsf M(\omega)\right]\widetilde{\vb x}
	=q\widetilde{\vb E},
\end{equation*}
故偶极矩 $\widetilde{\vb p}=q\widetilde{\vb x}=\alpha(\omega)\widetilde{\vb E}$ 的精确线性极化率为
\begin{equation*}
	\alpha(\omega)
	={q^2\over K-\omega^2\mathsf M(\omega)}.
\end{equation*}
定义物理质量 $m=\mathsf M(0)$ 和 $\omega_0^2=K/m$，并写
\begin{equation*}
	\Delta\mathsf M(\omega)=\mathsf M(\omega)-m.
\end{equation*}
则极化率中同一个分母改写为
\begin{equation*}
	K-\omega^2\mathsf M(\omega)
	=m(\omega_0^2-\omega^2)-\omega^2\Delta\mathsf M(\omega).
\end{equation*}
弱自作用时，极点 $\omega_{\mathrm p}=\omega_0+\delta\omega-i\Gamma/2$ 由分母改写式一阶给出
\begin{equation*}
	\delta\omega
	=-{\omega_0\over2m}\Re\Delta\mathsf M(\omega_0),
	\qquad
	\Gamma
	={\omega_0\over m}\Im\mathsf M(\omega_0).
\end{equation*}
所以频移来自同一响应的色散实部，线宽来自同一响应的耗散虚部。点粒子长波极限下
\begin{equation*}
	\Delta\mathsf M(\omega)=im\tau_{\mathrm{rr}}\omega,
\end{equation*}
故
\begin{equation*}
	\delta\omega=0,
	\qquad
	\Gamma=\tau_{\mathrm{rr}}\omega_0^2,
\end{equation*}
而极化率化为
\begin{equation*}
	\alpha(\omega)
	={q^2/m\over\omega_0^2-\omega^2-i\tau_{\mathrm{rr}}\omega^3}.
\end{equation*}
若还有非辐射宽度 $\Gamma'$，只需把分母再减去 $i\Gamma'\omega$。在共振附近可用 $\tau_{\mathrm{rr}}\omega^3\simeq\Gamma\omega$，总宽度为 $\Gamma_{\mathrm t}=\Gamma+\Gamma'$。

同一个极化率决定散射、消光和吸收：
\begin{equation*}
	\sigma_{\mathrm{sca}}(\omega)
	={\omega^4\over6\pi\varepsilon_0^2c^4}|\alpha(\omega)|^2,
	\qquad
	\sigma_{\mathrm{ext}}(\omega)
	={\omega\over\varepsilon_0c}\Im\alpha(\omega),
	\qquad
	\sigma_{\mathrm{abs}}=\sigma_{\mathrm{ext}}-\sigma_{\mathrm{sca}}.
\end{equation*}
当阻尼完全来自辐射时，将点极限极化率代入上式，可直接验证
\begin{equation*}
	\sigma_{\mathrm{ext}}=\sigma_{\mathrm{sca}},
	\qquad
	\sigma_{\mathrm{abs}}=0.
\end{equation*}
自由粒子极限 $\omega_0=0$ 且 $\omega\tau_{\mathrm{rr}}\ll1$ 时，同一个散射截面进一步化为 Thomson 值
\begin{equation*}
	\sigma_{\mathrm T}
	={q^4\over6\pi\varepsilon_0^2m^2c^4}
	={8\pi\over3}
	\left({q^2\over4\pi\varepsilon_0mc^2}\right)^2.
\end{equation*}
至此，附加惯性、辐射功率、线宽、频移和散射截面都由同一个因果响应 $\mathsf M(\omega)$ 贯通起来。

后面的二体库仑问题同时应用两条已经建立的主线：无反作用轨道上的频谱分析，以及由 Dirac 方程降阶得到的缓慢轨道反作用。质心系中的电偶极矩为
\begin{equation*}
	\vb p=q_D\vb r,
	\qquad
	q_D=\mu\left({q_1\over M_1}-{q_2\over M_2}\right).
\end{equation*}
在长波、弱反作用和降阶近似下，降阶后的自作用运动方程对相对坐标给出
\begin{equation*}
	\mu\ddot{\vb r}
	=\vb F(\vb r)
	+{q_D^2\over6\pi\varepsilon_0c^3}\dddot{\vb r}.
\end{equation*}
若 $q_D=0$，偶极频谱和偶极反作用同时消失，领先项转为电四极阶。因而库仑例题中先在无反作用轨道上求完整频谱，再把频谱积分得到的能量、角动量通量绝热地反馈到轨道常数，正是前面世界管守恒关系在弱辐射条件下的具体实现。

本节中三个常被写成"电磁质量"的系数现在可以按其推导位置区分。$U_{\mathrm{es}}/c^2$ 是静止系电磁能量除以 $c^2$；$4U_{\mathrm{es}}/(3c^2)$ 是球对称电磁子系统在实验室同时面上的零频平移惯性；$q^2/(8\pi\varepsilon_0c^2\epsilon)$ 是 Dirac 正交世界管中的奇异束缚动量系数。它们分别来自不同的、已经在上文明确写出的物理量，不能在省略源模型、积分超曲面和极限次序以后互相替换。

\subsubsection{例：有限质量二体库仑运动的辐射}

本例不再重新建立频谱理论，而是逐项实现前面的一般程序。首先在质心系中把两个可反冲点电荷的多极矩化为相对坐标的函数；随后求解统一包含吸引与排斥的库仑圆锥曲线。正相对能量的双曲轨道按照非周期过程的 Fourier 变换约定产生连续谱，负相对能量的椭圆轨道按照 Fourier 级数约定产生离散线谱；固定轨道的结果分别代入一般电偶极、电四极频域频谱公式，而散射系综则使用能量加权的一般谱辐射截面公式。这样，轨道动力学、特殊函数计算和辐射学归一化在逻辑上彼此分离，又能在最终结果中相互校验。

\textbf{质心系约化与统一轨道。}

设两粒子的质量和有符号电荷分别为 $(M_1,q_1)$、$(M_2,q_2)$。取
\begin{equation*}
	M=M_1+M_2,
	\qquad
	\mu=\frac{M_1M_2}{M},
	\qquad
	\vb R=\frac{M_1\vb r_1+M_2\vb r_2}{M},
	\qquad
	\vb r=\vb r_1-\vb r_2.
\end{equation*}
在质心系 $\vb R=0$ 中
\begin{equation*}
	\vb r_1=\frac{M_2}{M}\vb r,
	\qquad
	\vb r_2=-\frac{M_1}{M}\vb r.
\end{equation*}
因此电偶极矩约化为
\begin{equation*}
	\vb p=q_D\vb r,
	\qquad
	q_D=\frac{q_1M_2-q_2M_1}{M}
	=\mu\left(\frac{q_1}{M_1}-\frac{q_2}{M_2}\right),
\end{equation*}
无迹电四极矩约化为
\begin{equation*}
	D_{ij}=q_Q\left(3r_ir_j-r^2\delta_{ij}\right),
	\qquad
	q_Q=\frac{q_1M_2^2+q_2M_1^2}{M^2}
	=\mu^2\left(\frac{q_1}{M_1^2}+\frac{q_2}{M_2^2}\right).
\end{equation*}
轨道磁偶极矩则为
\begin{equation*}
	\vb m_{\mathrm{orb}}
	=\frac{1}{2}\sum_{a=1}^2q_a\vb r_a\times\dot{\vb r}_a
	=\frac{q_Q}{2\mu}\vb L,
	\qquad
	\vb L=\mu\vb r\times\dot{\vb r}.
\end{equation*}
对任意中心力零阶轨道，$\vb L$ 守恒，故 $\ddot{\vb m}_{\mathrm{orb}}=0$。于是本问题的正频率轨道辐射没有 $M1$ 项；保留至相对阶 $v^2/c^2$ 时，只需计算 $E1$ 与 $E2$。

库仑相互作用写成
\begin{equation*}
	V(r)=\sigma\frac{\mathcal C}{r},
	\qquad
	\mathcal C=\frac{|q_1q_2|}{4\pi\varepsilon_0},
	\qquad
	\sigma=\operatorname{sgn}(q_1q_2).
\end{equation*}
其中 $\sigma=+1$ 表示排斥，$\sigma=-1$ 表示吸引。相对运动拉格朗日量为
\begin{equation*}
	L_{\mathrm{rel}}=\frac{1}{2}\mu\dot{\vb r}^{\,2}-\sigma\frac{\mathcal C}{r},
\end{equation*}
从而
\begin{equation*}
	\mu\ddot{\vb r}=\sigma\mathcal C\frac{\vb r}{r^3}.
\end{equation*}
有限靶质量只通过 $\mu,q_D,q_Q$ 进入结果，并不破坏相对运动的中心力可积性。

在完成源的约化以后，考虑正相对能量的非束缚运动。设无穷远相对速度为 $v$，碰撞参数为 $b$，则
\begin{equation*}
	\mathcal E=\frac{1}{2}\mu v^2,
	\qquad
	L=\mu vb.
\end{equation*}
令 $w(\theta)=1/r$。由 Binet 方程
\begin{equation*}
	\frac{\d^2w}{\d\theta^2}+w
	=-\frac{\mu}{L^2w^2}\left(\sigma\mathcal Cw^2\right)
	=-\sigma\frac{\mu\mathcal C}{L^2},
\end{equation*}
并取近心点为 $\theta=0$，可得
\begin{equation*}
	r(\theta)=\frac{p}{\varepsilon\cos\theta-\sigma},
	\qquad
	p=\frac{L^2}{\mu\mathcal C},
	\qquad
	\varepsilon^2=1+\frac{2\mathcal EL^2}{\mu\mathcal C^2}.
\end{equation*}
定义无量纲参数
\begin{equation*}
	b_0=\frac{\mathcal C}{\mu v^2},
	\qquad
	\omega_0=\frac{v}{b_0},
	\qquad
	s=\frac{b}{b_0},
	\qquad
	\varepsilon=\sqrt{1+s^2},
	\qquad
	\nu=\frac{\omega}{\omega_0}.
\end{equation*}
此时 $p=b_0s^2$。为进行傅里叶积分，引入双曲参数 $u$：
\begin{equation*}
	\begin{aligned}
	x(u)&=b_0\left(\varepsilon+\sigma\cosh u\right),\\
	y(u)&=b_0s\sinh u,\\
	r(u)&=b_0\left(\varepsilon\cosh u+\sigma\right).
	\end{aligned}
\end{equation*}
确有 $x^2+y^2=r^2$，且由 $p=\varepsilon x-\sigma r$ 可恢复上面的圆锥曲线轨道。面积速度守恒给出
\begin{equation*}
	x\dot y-y\dot x=\frac{L}{\mu}=vb.
\end{equation*}
另一方面
\begin{equation*}
	x\frac{\d y}{\d u}-y\frac{\d x}{\d u}
	=b_0^2s\left(\varepsilon\cosh u+\sigma\right),
\end{equation*}
故
\begin{equation*}
	\frac{\d t}{\d u}
	=\frac{1}{\omega_0}\left(\varepsilon\cosh u+\sigma\right).
\end{equation*}
取近心点 $u=0$ 时 $t=0$，得到统一的时间参数式
\begin{equation*}
	t(u)=\frac{1}{\omega_0}\left(\varepsilon\sinh u+\sigma u\right).
\end{equation*}
空间参数式与时间参数式合并即为完整的双曲轨道参数化。最近距离为
\begin{equation*}
	r_{\min}=b_0(\varepsilon+\sigma),
\end{equation*}
散射角的绝对值满足
\begin{equation*}
	\tan\frac{\Theta}{2}=\frac{1}{s},
	\qquad
	\sin\frac{\Theta}{2}=\frac{1}{\varepsilon}.
\end{equation*}

\textbf{母积分与虚阶 Bessel 函数。}

为计算双曲轨道上的连续谱，需要先建立一个能够生成全部轨道积分的母公式。

所有连续谱积分均包含相位
\begin{equation*}
	\omega t(u)=\nu\left(\varepsilon\sinh u+\sigma u\right).
\end{equation*}
为统一处理 $\sinh u$、$\cosh u$ 的有限次幂，定义整数 $n$ 的母积分
\begin{equation*}
	\mathcal J_n^{(\sigma)}(\nu,\varepsilon)
	=\int_{-\infty}^{\infty}
	e^{nu}e^{i\nu(\varepsilon\sinh u+\sigma u)}\d u.
\end{equation*}
严格地说，可先在被积函数中加入收敛因子，再在积分完成后令收敛参数趋于零。将积分路径平移为 $u=w+i\pi/2$，利用
\begin{equation*}
	\sinh\left(w+\frac{i\pi}{2}\right)=i\cosh w
\end{equation*}
以及第二类修正 Bessel 函数的积分表示
\begin{equation*}
	K_\lambda(z)=\frac{1}{2}\int_{-\infty}^{\infty}
	e^{-z\cosh w+\lambda w}\d w,
	\qquad \Re z>0,
\end{equation*}
得到
\begin{equation*}
	\mathcal J_n^{(\sigma)}
	=2e^{-\sigma\pi\nu/2}e^{in\pi/2}
	K_{n+i\sigma\nu}(\nu\varepsilon).
\end{equation*}
特别地
\begin{equation*}
	\mathcal J_0^{(\sigma)}
	=2e^{-\sigma\pi\nu/2}K_{i\nu}(\nu\varepsilon),
\end{equation*}
因为 $K_{-i\nu}=K_{i\nu}$。以下简记
\begin{equation*}
	K=K_{i\nu}(\nu\varepsilon),
	\qquad
	K'=\left.\frac{\partial K_{i\nu}(x)}{\partial x}\right|_{x=\nu\varepsilon}.
\end{equation*}
整数平移阶数可由递推式
\begin{equation*}
	K_{\lambda+1}(x)=\frac{\lambda}{x}K_\lambda(x)-K_\lambda'(x),
	\qquad
	K_{\lambda-1}(x)=-\frac{\lambda}{x}K_\lambda(x)-K_\lambda'(x)
\end{equation*}
逐级约化，而二阶导数可由虚阶 Bessel 方程
\begin{equation*}
	x^2K_{i\nu}''(x)+xK_{i\nu}'(x)-(x^2-\nu^2)K_{i\nu}(x)=0
\end{equation*}
消去。因此有限次多项式乘以双曲轨道相位的积分最终都可写成 $K_{i\nu}$ 与 $K_{i\nu}'$ 的组合。

\textbf{电偶极与电四极连续谱。}

先计算电偶极项。由 $\vb p=q_D\vb r$，只需求相对加速度的傅里叶变换。从轨道参数式得到
\begin{equation*}
	\dot x=\sigma v\frac{\sinh u}{\varepsilon\cosh u+\sigma},
	\qquad
	\dot y=vs\frac{\cosh u}{\varepsilon\cosh u+\sigma}.
\end{equation*}
又有
\begin{equation*}
	\frac{\d}{\d u}e^{i\omega t(u)}
	=i\nu\left(\varepsilon\cosh u+\sigma\right)e^{i\omega t(u)}.
\end{equation*}
于是，在加入收敛因子后分部积分，
\begin{equation*}
	\begin{aligned}
	\widetilde a_x
	&=\int_{-\infty}^{\infty}\frac{\d\dot x}{\d u}e^{i\omega t}\d u
	=-i\sigma v\nu\int_{-\infty}^{\infty}\sinh u\,e^{i\omega t}\d u,\\
	\widetilde a_y
	&=\int_{-\infty}^{\infty}\frac{\d\dot y}{\d u}e^{i\omega t}\d u
	=-ivs\nu\int_{-\infty}^{\infty}\cosh u\,e^{i\omega t}\d u.
	\end{aligned}
\end{equation*}
第一项由
\begin{equation*}
	\frac{\partial\mathcal J_0^{(\sigma)}}{\partial\varepsilon}
	=i\nu\int_{-\infty}^{\infty}\sinh u\,e^{i\omega t}\d u
\end{equation*}
计算。第二项利用全导数积分
\begin{equation*}
	0=\int_{-\infty}^{\infty}\frac{\d}{\d u}e^{i\omega t}\d u
	=i\nu\int_{-\infty}^{\infty}
	\left(\varepsilon\cosh u+\sigma\right)e^{i\omega t}\d u
\end{equation*}
化为 $\mathcal J_0^{(\sigma)}$。代入母积分结果后得到
\begin{equation*}
	\widetilde a_x=-2\sigma v\nu e^{-\sigma\pi\nu/2}K',
	\qquad
	\widetilde a_y=2i\sigma v\nu\frac{s}{\varepsilon}e^{-\sigma\pi\nu/2}K.
\end{equation*}
因此
\begin{equation*}
	\widetilde{\ddot p}_i\widetilde{\ddot p}_i^*
	=4q_D^2v^2\nu^2e^{-\sigma\pi\nu}
	\left(K'^2+\frac{s^2}{\varepsilon^2}K^2\right),
\end{equation*}
代入一般电偶极频域频谱公式，得到任意质量比、任意碰撞参数的电偶极连续谱闭式
\begin{equation*}
	\frac{\d E_{E1}^{(\sigma)}(b)}{\d\omega}
	=\frac{2q_D^2v^2}{3\pi^2\varepsilon_0c^3}
	\nu^2e^{-\sigma\pi\nu}
	\left[
	K'^2+\frac{s^2}{\varepsilon^2}K^2
	\right].
\end{equation*}
吸引与排斥的差别全部包含在 $e^{-\sigma\pi\nu}$ 中；其余函数只依赖于 $\nu$ 与 $s=b/b_0$。

电四极项需要相对坐标二次型的傅里叶变换。由有效四极电荷，只需计算三个独立的二次轨道矩
\begin{equation*}
	R_{xx}=\int_{-\infty}^{\infty}x^2e^{i\omega t}\d t,
	\qquad
	R_{yy}=\int_{-\infty}^{\infty}y^2e^{i\omega t}\d t,
	\qquad
	R_{xy}=\int_{-\infty}^{\infty}xye^{i\omega t}\d t.
\end{equation*}
由双曲空间参数式及 $\d t/\d u=(\varepsilon\cosh u+\sigma)/\omega_0$，先写成
\begin{equation*}
	\begin{aligned}
	R_{xx}&=\frac{b_0^3}{v}\int_{-\infty}^{\infty}
	(\varepsilon+\sigma\cosh u)^2(\varepsilon\cosh u+\sigma)e^{i\omega t}\d u,\\
	R_{yy}&=\frac{b_0^3s^2}{v}\int_{-\infty}^{\infty}
	\sinh^2u(\varepsilon\cosh u+\sigma)e^{i\omega t}\d u,\\
	R_{xy}&=\frac{b_0^3s}{v}\int_{-\infty}^{\infty}
	\sinh u(\varepsilon+\sigma\cosh u)(\varepsilon\cosh u+\sigma)e^{i\omega t}\d u.
	\end{aligned}
\end{equation*}
为把每一项直接归入母积分，定义
\begin{equation*}
	\mathcal C_m=\frac{\mathcal J_m^{(\sigma)}+\mathcal J_{-m}^{(\sigma)}}{2},
	\qquad
	\mathcal S_m=\frac{\mathcal J_m^{(\sigma)}-\mathcal J_{-m}^{(\sigma)}}{2}.
\end{equation*}
其中 $\mathcal C_m$、$\mathcal S_m$ 分别是 $\cosh(mu)$、$\sinh(mu)$ 与轨道相位乘积的积分。利用
\begin{equation*}
	\cosh^2u=\frac{1+\cosh2u}{2},
	\quad
	\cosh^3u=\frac{3\cosh u+\cosh3u}{4},
	\quad
	\sinh u\cosh u=\frac{\sinh2u}{2},
\end{equation*}
\begin{equation*}
	\sinh u(\cosh^2u+1)=\frac{\sinh3u+5\sinh u}{4},
\end{equation*}
二次轨道矩积分逐项展开为
\begin{equation*}
	\begin{aligned}
	R_{xx}=\frac{b_0^3}{v}\bigg[&
	\frac{\varepsilon}{4}\mathcal C_3
	+\varepsilon\left(\varepsilon^2+\frac{11}{4}\right)\mathcal C_1
	+\frac{\sigma(2\varepsilon^2+1)}{2}\mathcal C_2
	+\frac{\sigma(4\varepsilon^2+1)}{2}\mathcal C_0\bigg],\\
	R_{yy}=\frac{b_0^3s^2}{v}\bigg[&
	\frac{\varepsilon}{4}\mathcal C_3
	-\frac{\varepsilon}{4}\mathcal C_1
	+\frac{\sigma}{2}\mathcal C_2
	-\frac{\sigma}{2}\mathcal C_0\bigg],\\
	R_{xy}=\frac{b_0^3s}{v}\bigg[&
	\frac{\varepsilon^2+1}{2}\mathcal S_2
	+\frac{\sigma\varepsilon}{4}
	\left(\mathcal S_3+5\mathcal S_1\right)\bigg].
	\end{aligned}
\end{equation*}
母积分结果将 $\mathcal C_m,\mathcal S_m$ 写成 $K_{i\sigma\nu\pm m}(\nu\varepsilon)$；再反复使用递推式，所有整数平移阶数消去。定义
\begin{equation*}
	\mathcal A_\sigma=\frac{4b_0^3}{v\nu^2}e^{-\sigma\pi\nu/2},
\end{equation*}
可得
\begin{equation*}
	\begin{aligned}
	R_{xx}&=\sigma\mathcal A_\sigma
	\left(\frac{\nu s^2}{\varepsilon}K'+\frac{1}{\varepsilon^2}K\right),\\
	R_{yy}&=-\sigma\mathcal A_\sigma
	\left(\frac{\nu s^2}{\varepsilon}K'-\frac{s^2}{\varepsilon^2}K\right),\\
	R_{xy}&=i\sigma\mathcal A_\sigma s
	\left(\frac{1}{\varepsilon}K'-\frac{\nu s^2}{\varepsilon^2}K\right).
	\end{aligned}
\end{equation*}
轨道位于 $xy$ 平面，故
\begin{equation*}
	\begin{aligned}
	\widetilde D_{xx}&=q_Q(2R_{xx}-R_{yy}),
	&\widetilde D_{yy}&=q_Q(2R_{yy}-R_{xx}),\\
	\widetilde D_{zz}&=-q_Q(R_{xx}+R_{yy}),
	&\widetilde D_{xy}&=\widetilde D_{yx}=3q_QR_{xy}.
	\end{aligned}
\end{equation*}
直接收缩得到
\begin{equation*}
	\widetilde D_{ij}\widetilde D_{ij}^*
	=6q_Q^2\left[
	|R_{xx}|^2+|R_{yy}|^2+3|R_{xy}|^2
	-\operatorname{Re}(R_{xx}R_{yy}^*)
	\right].
\end{equation*}
将二次轨道矩分量代入并整理，可写成
\begin{equation*}
	\widetilde D_{ij}\widetilde D_{ij}^*
	=\frac{96q_Q^2b_0^6}{v^2\nu^4}e^{-\sigma\pi\nu}\mathcal F_Q(\nu,s),
\end{equation*}
其中
\begin{equation*}
	\begin{aligned}
	\mathcal F_Q(\nu,s)=\frac{1}{\varepsilon^4}\bigg[&
	3\varepsilon^2s^2(1+\nu^2s^2)K'^2
	+3\varepsilon\nu s^2(1-3s^2)KK'\\
	&+\left(1-s^2+s^4+3\nu^2s^6\right)K^2
	\bigg].
	\end{aligned}
\end{equation*}
利用 $\omega=\nu v/b_0$，把收缩结果代入一般电四极频域频谱公式，得到电四极连续谱闭式
\begin{equation*}
	\frac{\d E_{E2}^{(\sigma)}(b)}{\d\omega}
	=\frac{2q_Q^2v^4}{15\pi^2\varepsilon_0c^5}
	\nu^2e^{-\sigma\pi\nu}\mathcal F_Q(\nu,s).
\end{equation*}
至此，有限质量二体库仑散射在固定碰撞参数下的 $E1$、$E2$ 连续谱均已化为虚阶第二类修正 Bessel 函数的闭式。
\textbf{总辐射能量与角动量。}

若不分辨频率，则同一结果可由 Parseval 恒等式写成
\begin{equation*}
	\Delta E_X^{(\sigma)}(b)
	=\int_0^\infty\frac{\d E_X^{(\sigma)}(b)}{\d\omega}\d\omega,
	\qquad X=E1,E2.
\end{equation*}
直接在时域积分可以更简洁地求出这个频率积分，并同时得到辐射角动量。令
\begin{equation*}
	f(\theta)=\varepsilon\cos\theta-\sigma,
	\qquad
	\theta_\infty=\arccos\frac{\sigma}{\varepsilon},
	\qquad
	s=\sqrt{\varepsilon^2-1},
\end{equation*}
则双曲轨道可写成
\begin{equation*}
	r=\frac{p}{f(\theta)},
	\qquad
	\frac{\d t}{\d\theta}
	=\sqrt{\frac{\mu p^3}{\mathcal C}}\frac{1}{f(\theta)^2},
	\qquad
	-\theta_\infty<\theta<\theta_\infty.
\end{equation*}
由 $\vb p=q_D\vb r$ 和统一库仑相对运动方程，
\begin{equation*}
	\ddot{\vb p}
	=\sigma\frac{q_D\mathcal C}{\mu}\frac{\vb r}{r^3},
	\qquad
	\dot{\vb p}\times\ddot{\vb p}
	=-\sigma\frac{q_D^2\mathcal C}{\mu^2r^3}\vb L.
\end{equation*}
代入前文的瞬时 $E1$ 能量与角动量通量，只需计算
\begin{equation*}
	\begin{aligned}
	\int_{-\theta_\infty}^{\theta_\infty}f(\theta)^2\d\theta
	&=(\varepsilon^2+2)\theta_\infty-3\sigma s,\\
	\int_{-\theta_\infty}^{\theta_\infty}f(\theta)\d\theta
	&=2\left(s-\sigma\theta_\infty\right).
	\end{aligned}
\end{equation*}
于是一次双曲散射的电偶极总辐射能量和辐射角动量分别为
\begin{equation*}
	\Delta E_{E1}^{(\sigma)}
	=\frac{q_D^2\mathcal C^{3/2}}
	{6\pi\varepsilon_0c^3\mu^{3/2}p^{5/2}}
	\left[(\varepsilon^2+2)\theta_\infty-3\sigma s\right],
\end{equation*}
\begin{equation*}
	\Delta\vb L_{E1,\mathrm{rad}}^{(\sigma)}
	=\frac{q_D^2\mathcal C}
	{3\pi\varepsilon_0c^3\mu p}
	\left(\theta_\infty-\sigma s\right)\uv L.
\end{equation*}
对于吸引散射，$\theta_\infty=\pi-\arccos(1/\varepsilon)$；对于排斥散射，$\theta_\infty=\arccos(1/\varepsilon)$。因此吸引与排斥两类双曲散射的偶极损失，都由上面两式给出。

四极总损失同样可从固定碰撞参数频谱积分得到，也可在时域直接完成。由
\begin{equation*}
	\begin{aligned}
	\dddot D_{ij}\dddot D_{ij}
	&=\frac{24q_Q^2\mathcal C^2}{\mu^2r^6}
	\left[12r^2\dot{\vb r}^{\,2}
	-11(\vb r\cdot\dot{\vb r})^2\right],\\
	\ddot{\vb D}\crossdot\dddot{\vb D}
	&=\frac{36\sigma\mathcal C q_Q^2}{\mu}
	(\vb r\times\dot{\vb r})
	\left[
	\frac{2\sigma\mathcal C}{\mu r^4}
	+\frac{3(\vb r\cdot\dot{\vb r})^2}{r^5}
	-\frac{2\dot{\vb r}^{\,2}}{r^3}
	\right],
	\end{aligned}
\end{equation*}
以及
\begin{equation*}
	\dot r=\frac{L\varepsilon\sin\theta}{\mu p},
	\qquad
	r\dot\theta=\frac{Lf(\theta)}{\mu p},
	\qquad
	L=\sqrt{\mu\mathcal C p},
\end{equation*}
四极能量积分归结为
\begin{equation*}
	\int_{-\theta_\infty}^{\theta_\infty}
	\left[12f^4+\varepsilon^2\sin^2\theta\,f^2\right]\d\theta
	=24\left[
	\left(1+\frac{73}{24}\varepsilon^2+\frac{37}{96}\varepsilon^4\right)\theta_\infty
	-\sigma s\frac{602+673\varepsilon^2}{288}
	\right],
\end{equation*}
而角动量积分为
\begin{equation*}
	\int_{-\theta_\infty}^{\theta_\infty}
	f\left[\varepsilon^2\sin^2\theta+2\sigma f-2f^2\right]\d\theta
	=\sigma(8+7\varepsilon^2)\theta_\infty
	-(13+2\varepsilon^2)s.
\end{equation*}
因此
\begin{equation*}
	\Delta E_{E2}^{(\sigma)}
	=\frac{4q_Q^2\mathcal C^{5/2}}
	{5\pi\varepsilon_0c^5\mu^{5/2}p^{7/2}}
	\left[
	\left(1+\frac{73}{24}\varepsilon^2+\frac{37}{96}\varepsilon^4\right)\theta_\infty
	-\sigma s\frac{602+673\varepsilon^2}{288}
	\right],
\end{equation*}
\begin{equation*}
	\Delta\vb L_{E2,\mathrm{rad}}^{(\sigma)}
	=\frac{q_Q^2\mathcal C^2}
	{10\pi\varepsilon_0c^5\mu^2p^2}
	\left[
	(8+7\varepsilon^2)\theta_\infty
	-\sigma(13+2\varepsilon^2)s
	\right]\uv L.
\end{equation*}
这些公式在 $q_D=0$ 时给出领先阶辐射；特别地，同号且荷质比相等的双曲散射正是上面两式取 $\sigma=+1$ 的情形。

\textbf{谱辐射截面。}

固定碰撞参数的谱给出单条经典轨道的辐射。按照能量加权的一般谱辐射截面定义，轴对称库仑散射的谱辐射截面为
\begin{equation*}
	\frac{\d\chi_X^{(\sigma)}}{\d\omega}
	=2\pi\int_0^\infty b\d b\,
	\frac{\d E_X^{(\sigma)}(b)}{\d\omega},
	\qquad X=E1,E2.
\end{equation*}
由于
\begin{equation*}
	b\d b=b_0^2\varepsilon\d\varepsilon,
	\qquad
	x=\nu\varepsilon,
\end{equation*}
偶极截面成为
\begin{equation*}
	\frac{\d\chi_{E1}^{(\sigma)}}{\d\omega}
	=\frac{4q_D^2v^2b_0^2}{3\pi\varepsilon_0c^3}e^{-\sigma\pi\nu}
	\int_\nu^\infty x\left[
	K_{i\nu}'(x)^2+
	\left(1-\frac{\nu^2}{x^2}\right)K_{i\nu}(x)^2
	\right]\d x.
\end{equation*}
由虚阶 Bessel 方程，
\begin{equation*}
	\frac{\d}{\d x}\left[xK_{i\nu}(x)K_{i\nu}'(x)\right]
	=x\left[
	K_{i\nu}'(x)^2+
	\left(1-\frac{\nu^2}{x^2}\right)K_{i\nu}(x)^2
	\right].
\end{equation*}
$K_{i\nu}(x)$ 在 $x\to\infty$ 指数衰减，故
\begin{equation*}
	\frac{\d\chi_{E1}^{(\sigma)}}{\d\omega}
	=\frac{4q_D^2v^2b_0^2}{3\pi\varepsilon_0c^3}
	e^{-\sigma\pi\nu}
	\left[-\nu K_{i\nu}(\nu)K_{i\nu}'(\nu)\right].
\end{equation*}
四极截面同样可以完全积出。把 $\mathcal F_Q$ 的定义中
\begin{equation*}
	\varepsilon=\frac{x}{\nu},
	\qquad
	s^2=\frac{x^2-\nu^2}{\nu^2}
\end{equation*}
代入，记 $K=K_{i\nu}(x)$，则可由虚阶 Bessel 方程逐项验证
\begin{equation*}
	x\mathcal F_Q
	=\frac{\d}{\d x}
	\left[A(x)K^2+B(x)KK'+C(x)K'^2\right],
\end{equation*}
其中
\begin{equation*}
	\begin{aligned}
	A(x)&=-\frac{(x^2-\nu^2)(2x^2-3\nu^2)}{2x^2},\\
	B(x)&=3\frac{(x^2-\nu^2)^2+x^2}{x},\\
	C(x)&=\frac{3\nu^2-13x^2}{2}.
	\end{aligned}
\end{equation*}
在 $x=\nu$ 处
\begin{equation*}
	A(\nu)=0,
	\qquad
	B(\nu)=3\nu,
	\qquad
	C(\nu)=-5\nu^2,
\end{equation*}
而上边界仍为零，故
\begin{equation*}
	\int_\nu^\infty x\mathcal F_Q\d x
	=5\nu^2K_{i\nu}'(\nu)^2
	-3\nu K_{i\nu}(\nu)K_{i\nu}'(\nu).
\end{equation*}
最终得到四极谱辐射截面的闭式
\begin{equation*}
	\frac{\d\chi_{E2}^{(\sigma)}}{\d\omega}
	=\frac{4q_Q^2v^4b_0^2}{15\pi\varepsilon_0c^5}
	e^{-\sigma\pi\nu}
	\left[
	5\nu^2K_{i\nu}'(\nu)^2
	-3\nu K_{i\nu}(\nu)K_{i\nu}'(\nu)
	\right].
\end{equation*}
若需要光子数截面，只需除以单光子能量：
\begin{equation*}
	\frac{\d\sigma_{\gamma,X}}{\d\omega}
	=\frac{1}{\hbar\omega}\frac{\d\chi_X}{\d\omega}.
\end{equation*}
\textbf{软频与硬频渐近。}

精确闭式建立后，其软频和硬频结构都可直接由虚阶 Bessel 函数的渐近式读出。

固定 $s$ 而令 $\nu\to0$ 时，
\begin{equation*}
	K_{i\nu}(\nu\varepsilon)
	=\ln\frac{2}{\nu\varepsilon}-\gamma_{\mathrm E}+o(1),
	\qquad
	K_{i\nu}'(\nu\varepsilon)
	=-\frac{1}{\nu\varepsilon}+o(\nu^{-1}).
\end{equation*}
代入固定碰撞参数的闭式，得到
\begin{equation*}
	\frac{\d E_{E1}}{\d\omega}
	\longrightarrow
	\frac{2q_D^2v^2}{3\pi^2\varepsilon_0c^3}
	\frac{1}{1+s^2}
	=\frac{q_D^2}{6\pi^2\varepsilon_0c^3}|\Delta\vb v|^2,
\end{equation*}
以及
\begin{equation*}
	\frac{\d E_{E2}}{\d\omega}
	\longrightarrow
	\frac{2q_Q^2v^4}{5\pi^2\varepsilon_0c^5}
	\frac{s^2}{(1+s^2)^2}
	=\frac{q_Q^2v^4}{10\pi^2\varepsilon_0c^5}\sin^2\Theta.
\end{equation*}
偶极软谱只取决于初末速度差。四极软谱在严格正碰时为零，这是因为初末速度虽反向，但张量 $v_iv_j$ 不变。

令
\begin{equation*}
	L_\nu=\ln\frac{2}{\nu}-\gamma_{\mathrm E},
\end{equation*}
由截面闭式可得
\begin{equation*}
	\frac{\d\chi_{E1}^{(\sigma)}}{\d\omega}
	=\frac{4q_D^2v^2b_0^2}{3\pi\varepsilon_0c^3}
	\left[L_\nu+o(1)\right],
\end{equation*}
\begin{equation*}
	\frac{\d\chi_{E2}^{(\sigma)}}{\d\omega}
	=\frac{4q_Q^2v^4b_0^2}{15\pi\varepsilon_0c^5}
	\left[5+3L_\nu+o(1)\right].
\end{equation*}
软频最低阶与 $\sigma$ 无关；对数来自越来越大的碰撞参数，实际体系中的屏蔽长度或有限观测时间将给出红外截断。

固定 $s>0$ 而令 $\nu\to\infty$，定义
\begin{equation*}
	\Delta(s)=s-\arctan s>0.
\end{equation*}
虚阶 Bessel 函数的陡降区渐近为
\begin{equation*}
	K_{i\nu}(\nu\varepsilon)
	\sim\sqrt{\frac{\pi}{2\nu s}}
	\exp\left[-\frac{\pi\nu}{2}-\nu\Delta(s)\right],
	\qquad
	K_{i\nu}'(\nu\varepsilon)
	\sim-\frac{s}{\varepsilon}K_{i\nu}(\nu\varepsilon).
\end{equation*}
故
\begin{equation*}
	\frac{\d E_{E1}^{(\sigma)}}{\d\omega}
	\sim
	\frac{2q_D^2v^2}{3\pi\varepsilon_0c^3}
	\frac{\nu s}{\varepsilon^2}
	\exp\left[-2\nu\Delta(s)-(1+\sigma)\pi\nu\right],
\end{equation*}
\begin{equation*}
	\frac{\d E_{E2}^{(\sigma)}}{\d\omega}
	\sim
	\frac{2q_Q^2v^4}{5\pi\varepsilon_0c^5}
	\frac{\nu^3s^5}{\varepsilon^4}
	\exp\left[-2\nu\Delta(s)-(1+\sigma)\pi\nu\right].
\end{equation*}
对吸引势，$s\ll1$ 时 $\Delta(s)=s^3/3+O(s^5)$，故固定碰撞参数谱的特征截止满足 $\nu s^3\sim1$。上述固定 $s$ 渐近对 $s\to0$ 不一致；谱截面的高频极限必须使用转折点 $x=\nu$ 附近的一致渐近式
\begin{equation*}
	e^{\pi\nu/2}K_{i\nu}(\nu)
	\sim\frac{\sqrt3\,\Gamma(1/3)}{6^{2/3}\nu^{1/3}},
	\qquad
	-e^{\pi\nu/2}K_{i\nu}'(\nu)
	\sim\frac{\sqrt3\,\Gamma(2/3)}{6^{1/3}\nu^{2/3}}.
\end{equation*}
代入截面闭式，得到
\begin{equation*}
	\frac{\d\chi_{E1}^{(\sigma)}}{\d\omega}
	\sim
	\frac{4q_D^2v^2b_0^2}{3\sqrt3\,\varepsilon_0c^3}
	e^{-(1+\sigma)\pi\nu},
\end{equation*}
以及
\begin{equation*}
	\frac{\d\chi_{E2}^{(\sigma)}}{\d\omega}
	\sim
	\frac{8\sqrt\pi}{3^{7/6}\Gamma(1/6)}
	\frac{q_Q^2v^4b_0^2}{\varepsilon_0c^5}
	\nu^{2/3}e^{-(1+\sigma)\pi\nu}.
\end{equation*}
因此排斥势具有额外的 $e^{-2\pi\nu}$ 截止；理想点吸引势的偶极截面趋于常数，而四极截面按 $\nu^{2/3}$ 增长。后两者来自越来越小的碰撞参数，只是经典点势模型的中间渐近，最终会被有限尺寸、相对论、量子反冲或辐射反作用截断。

\textbf{相对论固定中心库仑散射的软谱。}

前面的双曲轨道闭式建立在非相对论二体动力学上。若只要求软频极限，则可以放宽速度条件而不必求完整的相对论 Fourier 积分，因为一般相对论软谱公式只需要初末速度。为给出一个完全解析的例子，考虑质量为 $m$、电荷为 $q$ 的试探粒子在固定点电荷产生的外部库仑势中运动：
\begin{equation*}
	V_\sigma(r)=\sigma\frac{\kappa}{r},
	\qquad
	\kappa=\frac{|qQ|}{4\pi\varepsilon_0},
	\qquad
	\sigma=\operatorname{sgn}(qQ).
\end{equation*}
这里固定中心不辐射，全部软辐射来自试探粒子的速度改变。应当强调，这一模型是外部静电场中的相对论单粒子问题；若两个有限质量电荷都相对论性，完整 Maxwell--Lorentz 动力学不能简单化为一个瞬时的约化质量中心势。

试探粒子的守恒能量和角动量为
\begin{equation*}
	\mathcal E
	=\sqrt{m^2c^4+c^2\left(p_r^2+\frac{L^2}{r^2}\right)}
	+\sigma\frac{\kappa}{r},
	\qquad
	L=p_\infty b,
\end{equation*}
其中
\begin{equation*}
	\mathcal E=\gamma_\infty mc^2,
	\qquad
	p_\infty=\gamma_\infty mv_\infty,
	\qquad
	\beta_\infty=\frac{v_\infty}{c}.
\end{equation*}
令 $u=1/r$。由 Hamilton 方程可得 $p_r=-L\d u/\d\varphi$，于是径向能量方程化为
\begin{equation*}
	\left(\frac{\d u}{\d\varphi}\right)^2
	=\frac{p_\infty^2}{L^2}
	-\frac{2\sigma\mathcal E\kappa}{L^2c^2}u
	-\left(1-\frac{\kappa^2}{L^2c^2}\right)u^2.
\end{equation*}
对 $\varphi$ 求导，得到线性轨道方程
\begin{equation*}
	\frac{\d^2u}{\d\varphi^2}+\lambda^2u
	=-\frac{\sigma\mathcal E\kappa}{L^2c^2},
	\qquad
	\lambda=\sqrt{1-a^2},
	\qquad
	a=\frac{\kappa}{Lc}.
\end{equation*}
先考虑 $a<1$ 的散射支。把最近点取在 $\varphi=0$，解为
\begin{equation*}
	u(\varphi)=A\cos(\lambda\varphi)-\sigma u_c,
	\qquad
	u_c=\frac{\mathcal E\kappa}{L^2c^2-\kappa^2},
\end{equation*}
其中由径向能量方程的常数项得到
\begin{equation*}
	A^2=u_c^2+\frac{p_\infty^2}{L^2\lambda^2}.
\end{equation*}
令
\begin{equation*}
	\delta
	=\arctan\frac{a}{\beta_\infty\lambda}.
\end{equation*}
利用 $\mathcal E/(p_\infty c)=1/\beta_\infty$，可验证
\begin{equation*}
	\frac{u_c}{A}
	=\frac{a/(\beta_\infty\lambda)}
	{\sqrt{1+a^2/(\beta_\infty^2\lambda^2)}}
	=\sin\delta.
\end{equation*}
轨道的两条渐近线满足 $u(\pm\varphi_\infty)=0$，因而
\begin{equation*}
	\lambda\varphi_\infty=\frac{\pi}{2}-\sigma\delta.
\end{equation*}
从入射无穷远到出射无穷远所扫过的极角为
\begin{equation*}
	\Delta\varphi_\sigma
	=2\varphi_\infty
	=\frac{\pi-2\sigma\delta}{\lambda},
\end{equation*}
相对于无相互作用直线的有符号偏转角为
\begin{equation*}
	\chi_\sigma
	=\pi-\Delta\varphi_\sigma
	=\pi-\frac{\pi-2\sigma\arctan[a/(\beta_\infty\lambda)]}{\lambda}.
\end{equation*}
这就是相对论固定中心库仑散射的 Darwin 角。当轨道发生绕转时，$\chi_\sigma$ 可以超出通常的主值区间；软谱只通过初末速度的内积依赖它，因此应保留 $\cos\chi_\sigma$，而不必人为把角度折回 $[0,\pi]$。吸引势在 $a\geq1$ 时失去离心势垒并出现落向中心的轨道，上面的偏转角公式所描述的是 $a<1$ 的普通散射支。

初末速率相同，均为 $\beta_\infty c$，故两条渐近外线之间的相对 Lorentz 因子为
\begin{equation*}
	\Gamma_\sigma(b)
	=\gamma_\infty^2
	\left(1-\beta_\infty^2\cos\chi_\sigma(b)\right)
	=1+2\gamma_\infty^2\beta_\infty^2
	\sin^2\frac{\chi_\sigma(b)}{2}.
\end{equation*}
于是无需求解相对论轨道的时间参数或完整 Lienard--Wiechert Fourier 积分，便可由单电荷相对论软谱公式立即得到固定碰撞参数的解析软谱
\begin{equation*}
	\left.\frac{\d E_{\mathrm{soft}}^{(\sigma)}(b)}{\d\omega}\right|_{\omega\to0}
	=\frac{q^2}{2\pi^2\varepsilon_0c}
	\left[
	\frac{\Gamma_\sigma(b)\operatorname{arccosh}\Gamma_\sigma(b)}
	{\sqrt{\Gamma_\sigma(b)^2-1}}-1
	\right].
\end{equation*}
相应的角分布同样是闭式：
\begin{equation*}
	\left.\frac{\d^2E_{\mathrm{soft}}^{(\sigma)}}{\d\omega\d\Omega}\right|_{\omega\to0}
	=\frac{q^2}{16\pi^3\varepsilon_0c}
	\left|
	\vb n\times\left\{\vb n\times
	\left[
	\frac{\boldsymbol\beta_f}{1-\vb n\cdot\boldsymbol\beta_f}
	-\frac{\boldsymbol\beta_i}{1-\vb n\cdot\boldsymbol\beta_i}
	\right]\right\}
	\right|^2,
\end{equation*}
其中 $|\boldsymbol\beta_i|=|\boldsymbol\beta_f|=\beta_\infty$ 且 $\boldsymbol\beta_i\cdot\boldsymbol\beta_f=\beta_\infty^2\cos\chi_\sigma$。

几个极限可用来检验这一闭式。首先，当 $a\ll1$ 时，
\begin{equation*}
	\lambda=1+O(a^2),
	\qquad
	\chi_\sigma
	=\frac{2\sigma\kappa}{\gamma_\infty m v_\infty^2b}+O(\kappa^2).
\end{equation*}
若再取 $\beta_\infty\ll1$，软谱闭式退化为
\begin{equation*}
	\left.\frac{\d E_{\mathrm{soft}}}{\d\omega}\right|_{\omega\to0}
	=\frac{2q^2v_\infty^2}{3\pi^2\varepsilon_0c^3}
	\sin^2\frac{\chi_\sigma}{2}
	+O(c^{-5}),
\end{equation*}
并由非相对论 Rutherford 关系 $\sin^2(\chi/2)=1/(1+b^2/b_0^2)$ 恢复固定碰撞参数偶极软谱的固定中心极限。

其次，在大碰撞参数端，把小偏转角展开与小相对速度展开结合给出
\begin{equation*}
	\left.\frac{\d E_{\mathrm{soft}}(b)}{\d\omega}\right|_{b\to\infty}
	\simeq
	\frac{2q^2\kappa^2}{3\pi^2\varepsilon_0m^2v_\infty^2c^3}\frac{1}{b^2}.
\end{equation*}
因此未屏蔽库仑场的谱辐射截面仍有 Coulomb 对数：
\begin{equation*}
	\left.\frac{\d\chi_{\mathrm{soft}}}{\d\omega}\right|_{\omega\to0}
	\simeq
	\frac{4q^2\kappa^2}{3\pi\varepsilon_0m^2v_\infty^2c^3}
	\ln\frac{b_{\max}}{b_{\min}}.
\end{equation*}
相对论因子在这个大 $b$ 的领先对数系数中抵消；速度和强场修正进入有限项以及小碰撞参数区域。最后，当 $\gamma_\infty\sin(|\chi_\sigma|/2)\gg1$ 时，软谱闭式变为
\begin{equation*}
	\left.\frac{\d E_{\mathrm{soft}}}{\d\omega}\right|_{\omega\to0}
	\simeq
	\frac{q^2}{2\pi^2\varepsilon_0c}
	\left[
	2\ln\left(2\gamma_\infty\left|\sin\frac{\chi_\sigma}{2}\right|\right)-1
	\right],
\end{equation*}
显示了相对论软轫致辐射的对数增强和前向束射。

软谱闭式的精确性具有明确含义：它对 $\beta_\infty$ 和偏转角不作展开，对 $a<1$ 的固定中心相对论库仑散射轨道使用精确的 Darwin 角，但只给出 $\omega\to0$ 的领先项。其软条件是辐射周期远长于有效散射时间，即 $\omega\tau_{\mathrm{sc}}\ll1$；它不要求 $\beta_\infty\ll1$，也不要求电偶极长波条件 $\omega a_{\mathrm{eff}}/c\ll1$。若固定中心近似被放宽而靶粒子也发生反冲，一般的相对论软谱积分公式仍然成立，只需把两粒子的入射、出射四条外线全部代入；困难仅转移到相对论二体散射角和渐近四动量的动力学求解上。
\textbf{偶极抵消的特殊情形。}

有限质量二体问题还有一个重要的特殊情形，即电偶极矩由于荷质比关系而抵消。

当两粒子的荷质比相同，
\begin{equation*}
	\frac{q_1}{M_1}=\frac{q_2}{M_2},
\end{equation*}
有效偶极电荷公式给出 $q_D=0$，整个电偶极谱消失，而固定碰撞参数的电四极谱与其截面闭式成为领先结果。若 $q_D\neq0$，由软频结果得到
\begin{equation*}
	\frac{\d E_{E2}/\d\omega}{\d E_{E1}/\d\omega}
	\longrightarrow
	\frac{3}{5}\left(\frac{q_Q}{q_D}\right)^2
	\frac{v^2}{c^2}\frac{s^2}{1+s^2}.
\end{equation*}
因此在没有电荷质量组合抵消时，四极通常比偶极低相对阶 $v^2/c^2$；当 $|q_D|\lesssim(v/c)|q_Q|$ 时，二者才可能在同一量级竞争。相反，若 $M_1=M_2$ 且 $q_1=-q_2$，则 $q_Q=0$ 而 $q_D=q_1$，最低阶仍为纯电偶极辐射。

\textbf{椭圆束缚运动的谱线与绝热演化。}

负相对能量只可能出现在吸引势 $\sigma=-1$，此时轨道为周期性椭圆。按照周期运动的 Fourier 级数约定与一般周期谱线功率公式，连续 Fourier 变换应改为离散 Fourier 级数，辐射由整数谐波线组成。设椭圆半长轴为 $a$，离心率为 $0\leq e<1$，平均角频率为
\begin{equation*}
	\Omega=\sqrt{\frac{\mathcal C}{\mu a^3}}.
\end{equation*}
以偏近点角 $u$ 参数化
\begin{equation*}
	x=a(\cos u-e),
	\qquad
	y=a\sqrt{1-e^2}\sin u,
	\qquad
	\Omega t=u-e\sin u.
\end{equation*}
按照周期 Fourier 级数约定，需要计算
\begin{equation*}
	\mathcal I_{N,m}(e)
	=\frac{1}{2\pi}\int_0^{2\pi}
	e^{imu}(1-e\cos u)e^{iN(u-e\sin u)}\d u.
\end{equation*}
利用 Jacobi--Anger 展开
\begin{equation*}
	e^{-iNe\sin u}=\sum_{k=-\infty}^{\infty}J_k(-Ne)e^{iku}
\end{equation*}
并逐项积分，得到
\begin{equation*}
	\mathcal I_{N,m}(e)
	=J_{N+m}(Ne)-\frac{e}{2}
	\left[J_{N+m-1}(Ne)+J_{N+m+1}(Ne)\right].
\end{equation*}
对 $N\geq1$ 有 $\mathcal I_{N,0}=0$。定义
\begin{equation*}
	C_{N,m}=\frac{\mathcal I_{N,m}+\mathcal I_{N,-m}}{2},
	\qquad
	S_{N,m}=\frac{\mathcal I_{N,m}-\mathcal I_{N,-m}}{2i}.
\end{equation*}
坐标的 Fourier 系数为
\begin{equation*}
	r_{x,N}=\frac{a}{N}J_N'(Ne),
	\qquad
	r_{y,N}=i\frac{a\sqrt{1-e^2}}{Ne}J_N(Ne).
\end{equation*}
代入一般周期谱线功率公式，得到第 $N$ 条电偶极谱线
\begin{equation*}
	P_{E1,N}
	=\frac{q_D^2a^2\Omega^4}{3\pi\varepsilon_0c^3}N^2
	\left[J_N'(Ne)^2+\frac{1-e^2}{e^2}J_N(Ne)^2\right],
	\qquad N\geq1.
\end{equation*}
四极所需的二次轨道矩系数为
\begin{equation*}
	\begin{aligned}
	R_{xx,N}&=a^2\left(\frac{1}{2}C_{N,2}-2eC_{N,1}\right),\\
	R_{yy,N}&=-\frac{a^2(1-e^2)}{2}C_{N,2},\\
	R_{xy,N}&=a^2\sqrt{1-e^2}
	\left(\frac{1}{2}S_{N,2}-eS_{N,1}\right).
	\end{aligned}
\end{equation*}
相应的 $D_{ij,N}$ 仍由双曲情形中四极分量与二次轨道矩的关系给出，因此
\begin{equation*}
	P_{E2,N}
	=\frac{q_Q^2(N\Omega)^6}{60\pi\varepsilon_0c^5}
	\left[
	|R_{xx,N}|^2+|R_{yy,N}|^2+3|R_{xy,N}|^2
	-\operatorname{Re}(R_{xx,N}R_{yy,N}^*)
	\right].
\end{equation*}
总平均谱写成
\begin{equation*}
	\frac{\d\overline P}{\d\omega}
	=\sum_{N=1}^{\infty}
	\left(P_{E1,N}+P_{E2,N}+\cdots\right)
	\delta(\omega-N\Omega).
\end{equation*}
对所有谐波求和必须恢复时域中的周期平均功率。以电偶极项为例，令
\begin{equation*}
	p=a(1-e^2),
	\qquad
	r=\frac{p}{1+e\cos\theta},
	\qquad
	\frac{\d t}{\d\theta}
	=\sqrt{\frac{\mu p^3}{\mathcal C}}
	\frac{1}{(1+e\cos\theta)^2},
	\qquad
	T=\frac{2\pi}{\Omega},
\end{equation*}
则由瞬时多极通量或等价地由谱线 Parseval 求和得到
\begin{equation*}
	-\frac{\d E_{\mathrm{orb}}}{\d t}
	=\sum_{N=1}^{\infty}P_{E1,N}
	=\frac{q_D^2\mathcal C^2}
	{6\pi\varepsilon_0c^3\mu^2p^4}
	\left(1+\frac{e^2}{2}\right)(1-e^2)^{3/2},
\end{equation*}
以及
\begin{equation*}
	-\frac{\d\vb L}{\d t}
	=\frac{q_D^2\mathcal C^{3/2}}
	{6\pi\varepsilon_0c^3\mu^{3/2}p^{5/2}}
	(1-e^2)^{3/2}\uv L.
\end{equation*}
这里 $E_{\mathrm{orb}}=-\mathcal C(1-e^2)/(2p)$ 且 $L=\sqrt{\mu\mathcal C p}$。由上面两式消去 $E_{\mathrm{orb}}$ 与 $L$，得到缓慢辐射反作用下的轨道参数方程
\begin{equation*}
	\frac{\dot p}{p}
	=-\frac{q_D^2\mathcal C}{3\pi\varepsilon_0c^3\mu^2p^3}(1-e^2)^{3/2},
	\qquad
	\frac{\dot e}{e}
	=-\frac{q_D^2\mathcal C}{4\pi\varepsilon_0c^3\mu^2p^3}(1-e^2)^{3/2}.
\end{equation*}
因此
\begin{equation*}
	\frac{e}{e_i}=\left(\frac{p}{p_i}\right)^{3/4}.
\end{equation*}
若记
\begin{equation*}
	\tau=\frac{\pi\varepsilon_0c^3\mu^2p_i^3}{\mathcal Cq_D^2},
	\qquad
	F(e)=\sqrt{1-e^2}+\frac{1}{\sqrt{1-e^2}},
\end{equation*}
则 $e_i\neq0$ 时的积分解可以写成简洁的隐式形式
\begin{equation*}
	t=\frac{4\tau}{e_i^4}\left[F(e_i)-F(e)\right],
	\qquad
	p=p_i\left(\frac{e}{e_i}\right)^{4/3}.
\end{equation*}
在这一绝热模型中，形式上的并合时刻为
\begin{equation*}
	t_{\mathrm c}
	=\frac{4\tau}{e_i^4}\left[F(e_i)-2\right].
\end{equation*}
圆轨道极限 $e_i\to0$ 给出
\begin{equation*}
	e(t)=0,
	\qquad
	p(t)=p_i\left(1-\frac{t}{\tau}\right)^{1/3},
	\qquad
	t_{\mathrm c}=\tau.
\end{equation*}
接近形式并合时，局域速度、有限尺寸和辐射反作用均不再是小修正，因此隐式解只描述缓慢演化的非相对论阶段。

圆轨道 $e=0$ 时，偶极矩只含基频，四极矩只含二次谐波：
\begin{equation*}
	P_{E1,1}=\frac{q_D^2a^2\Omega^4}{6\pi\varepsilon_0c^3},
	\qquad
	P_{E2,2}=\frac{2q_Q^2a^4\Omega^6}{5\pi\varepsilon_0c^5},
\end{equation*}
其余谱线均为零。若再有 $q_D=0$，最低非零辐射线便位于 $2\Omega$。

\textbf{适用范围。}

最后说明上述闭式的适用范围。这些解析式对质量比和库仑偏转角没有作展开，但仍属于非相对论长波多极理论。首先必须要求整个有效轨道上的局域相对速度远小于 $c$。排斥轨道上速度的最大值是无穷远速度，故 $v/c\ll1$ 基本足够；吸引轨道的近心速度为
\begin{equation*}
	v_{\mathrm p}=v\sqrt{\frac{\varepsilon+1}{\varepsilon-1}},
\end{equation*}
因而还需 $v_{\mathrm p}/c\ll1$。其次，多极展开要求对产生给定频率的有效空间尺度有 $\omega r_{\mathrm{eff}}/c\ll1$。排斥硬频区相应为 $\beta\nu\ll1$；吸引硬频截面的主导碰撞满足 $s\sim\nu^{-1/3}$，相应条件为 $\beta\nu^{1/3}\ll1$。

此外还应满足辐射反作用所造成的能量损失远小于入射相对动能，经典轨道角动量 $L/\hbar=\mu vb/\hbar$ 足够大，光子反冲满足 $\hbar\omega\ll\mu v^2/2$，并且主要贡献的最近距离大于粒子或核的有限尺寸、主要贡献的最大距离小于屏蔽尺度。在这些条件共同成立的区域，固定碰撞参数的 $E1$、$E2$ 连续谱闭式与相应的截面闭式给出非束缚二体库仑运动的连续谱；$E1$、$E2$ 谱线功率公式给出吸引束缚轨道的离散谱。

\begin{notes}
作为类比，广义相对论中最低阶引力波辐射由质量四极矩贡献。考虑做二体椭圆运动的质量为 \(M_1\) 和 \(M_2\) 的质点，记
\begin{equation*}
	\alpha = G M_1 M_2, \quad \mu = \frac{M_1 M_2}{M_1 + M_2}.
\end{equation*}

设 2 指向 1 的位矢为 \(\vb{r}\)，则二体运动方程、二体角动量、质心系中体系的质量四极矩为
\begin{equation*}
	\mu \frac{\d^2\vb{r}}{\d t^2} = -\frac{\alpha}{r^3}\vb{r}, \quad \vb{L} = \mu \vb{r} \times \frac{\d \vb{r}}{\d t}, \quad \vb{Q} = \mu \left( 3\vb{r}\vb{r} - r^2 \vb{I} \right).
\end{equation*}

进而，质量四极矩的各阶变化率为
\begin{equation*}
	\begin{aligned}
	\frac{\d \vb{Q}}{\d t} &= \mu \left( 3\frac{\d \vb{r}}{\d t}\vb{r} + 3\vb{r}\frac{\d \vb{r}}{\d t} - 2\left( \vb{r} \cdot \frac{\d \vb{r}}{\d t} \right)\vb{I} \right), \\
	\frac{\d^2\vb{Q}}{\d t^2} &= 2\alpha p \left( -\frac{3\vb{r}\vb{r} - r^2 \vb{I}}{r^3 p} + \frac{\mu^2}{L^2} \left( 3\frac{\d \vb{r}}{\d t}\frac{\d \vb{r}}{\d t} - \left| \frac{\d \vb{r}}{\d t} \right|^2 \vb{I} \right) \right), \\
	\frac{\d^3\vb{Q}}{\d t^3} &= \frac{2\alpha}{r^3} \left( 9\frac{\vb{r}\vb{r}}{r^2} \left( \vb{r} \cdot \frac{\d \vb{r}}{\d t} \right) - 6\left( \frac{\d \vb{r}}{\d t}\vb{r} + \vb{r}\frac{\d \vb{r}}{\d t} \right) + \left( \vb{r} \cdot \frac{\d \vb{r}}{\d t} \right)\vb{I} \right).\\
	\frac{\d^3\vb{Q}}{\d t^3} : \frac{\d^3\vb{Q}}{\d t^3} &= \frac{24\alpha^2}{r^6} \left( -11\left( \vb{r} \cdot \frac{\d \vb{r}}{\d t} \right)^2 + 12 r^2 \left| \frac{\d \vb{r}}{\d t} \right|^2 \right).\\
	\frac{\d^2\vb{Q}}{\d t^2} \crossdot \left( \frac{\d^3\vb{Q}}{\d t^3} \right) &= \frac{36\alpha^2 p \vb{L}}{\mu r^5} \left( \frac{2r}{p} - \frac{3\mu^2}{L^2} \left( \vb{r} \cdot \frac{\d \vb{r}}{\d t} \right)^2 + \frac{2\mu^2}{L^2} r^2 \left| \frac{\d \vb{r}}{\d t} \right|^2 \right).
	\end{aligned}
\end{equation*}

由于观测时更关心周期性椭圆运动的平均效应，因此考虑周期平均后的辐射损失能量、角动量。运用开普勒运动的结论
\begin{equation*}
	r = \frac{p}{1 + e\cos\theta}, \quad \frac{\d \vb{r}}{\d t} = \frac{L}{\mu r} \left( \frac{e\sin\theta}{1 + e\cos\theta}\uv{r} + \uv{\theta} \right), \quad L = \sqrt{\mu\alpha p}, \quad T = \frac{2\pi}{(1 - e^2)^{3/2}} \sqrt{\frac{p^3}{\alpha\mu}}.
\end{equation*}

并代入
\begin{equation*}
	\frac{\d t}{\d \theta} = \sqrt{\frac{\mu p^3}{\alpha}} \frac{1}{(1 + e\cos\theta)^2}.
\end{equation*}

经过繁琐的积分得
\begin{equation*}
	\begin{aligned}
	-\frac{\d E}{\d t} &= \frac{G}{45 c^5 T} \int_{0}^{2\pi} \left( \frac{\d^3\vb{Q}}{\d t^3} : \frac{\d^3\vb{Q}}{\d t^3} \right) \frac{\d t}{\d \theta} \d \theta
	= \frac{32 G}{5 c^5} \frac{\alpha^3}{\mu p^5} (1 - e^2)^{3/2} \left( 1 + \frac{73}{24}e^2 + \frac{37}{96}e^4 \right), \\
	-\frac{\d \vb{L}}{\d t} &= \frac{2 G}{45 c^5 T} \int_{0}^{2\pi} \left( \frac{\d^2\vb{Q}}{\d t^2} \crossdot \left( \frac{\d^3\vb{Q}}{\d t^3} \right) \right) \frac{\d t}{\d \theta} \d \theta
	= \frac{32 G}{5 c^5} \frac{\alpha^{5/2} \uv{L}}{\mu^{1/2} p^{7/2}} (1 - e^2)^{3/2} \left( 1 + \frac{7}{8}e^2 \right).
	\end{aligned}
\end{equation*}

用两质点质量及轨道半长轴 \(a = p(1-e^2)^{-1}\) 表示，即为
\begin{equation*}
	\begin{aligned}
	-\frac{\d E}{\d t} &= \frac{32 G^4 M_1^2 M_2^2 (M_1 + M_2)}{5 c^5 a^5} (1 - e^2)^{-7/2} \left( 1 + \frac{73}{24}e^2 + \frac{37}{96}e^4 \right), \\
	-\frac{\d \vb{L}}{\d t} &= \frac{32 G^{7/2} M_1^2 M_2^2 (M_1 + M_2)^{1/2}}{5 c^5 a^{7/2}} (1 - e^2)^{-2} \left( 1 + \frac{7}{8}e^2 \right).
	\end{aligned}
\end{equation*}
\end{notes}
